\documentclass[14pt,a4paper]{extarticle}
\usepackage[utf8]{inputenc}
\usepackage[T2A]{fontenc}
\usepackage[russian]{babel}
\usepackage{geometry}
\usepackage{setspace}

\geometry{left=25mm, right=15mm, top=20mm, bottom=20mm}
\onehalfspacing

\begin{document}

\begin{center}
    \Large \textbf{Речь для защиты презентации (Шпаргалка)} \\
    \vspace{0.3cm}
    \large \textit{Тема: Стандарты ЕСПД и Формализация задачи (Вариант 5)} \\
    \vspace{0.3cm}
    \normalsize Студент: Егоров А.В. (М3О-225БВ-24)
\end{center}
\vspace{1cm}

\textbf{[Слайд 1: Титульный]} \\
Здравствуйте, уважаемый Михаил Николаевич и коллеги! Меня зовут Егоров Александр. Мой сегодняшний доклад в рамках технологической практики состоит из двух частей. Первая часть посвящена стандартам Единой системы программной документации (ЕСПД), а вторая — формализации задачи и математическим расчетам по моему пятому варианту.

\vspace{0.5cm}
\textbf{[Слайд 2: Содержание]} \\
\textit{(Можно просто переключить на следующий слайд, не комментируя)}

\vspace{0.5cm}
\textbf{[Слайд 3: ГОСТ 19.502-78. Назначение]} \\
Начнем с ГОСТ 19.502-78, который регламентирует документ «Описание применения». Этот документ создается не для программистов, а исключительно для специалистов, эксплуатирующих программу. Его главная цель — дать исчерпывающие сведения о том, как внедрить программу и в каких условиях она гарантированно будет работать.

\vspace{0.5cm}
\textbf{[Слайд 4: ГОСТ 19.502-78. Состав]} \\
Структура документа строго регламентирована и включает:
1. Назначение программы (возможности и ограничения).
2. Условия применения (системные требования к железу и ОС).
3. Описание задачи (формализация и методы решения).
4. Входные и выходные данные (форматы файлов или пользовательского ввода).

\vspace{0.5cm}
\textbf{[Слайд 5: ГОСТ 19.005-85. Назначение]} \\
Второй стандарт — ГОСТ 19.005-85 — описывает правила выполнения схем алгоритмов и программ. Он базируется на международном стандарте ISO 5807-85, что позволяет нашим блок-схемам быть понятными инженерам по всему миру.

\vspace{0.5cm}
\textbf{[Слайд 6: ГОСТ 19.005-85. Оформление]} \\
Стандарт жестко описывает геометрию чертежа: линии потока должны быть строго вертикальными или горизонтальными. Стрелки ставятся только если поток идет нестандартно (снизу вверх или справа налево). А геометрия блоков привязана к сетке с шагом 5 мм (высота блока составляет половину или две трети от ширины).

\vspace{0.5cm}
\textbf{[Слайд 7: Постановка задачи (Вариант 5)]} \\
Теперь перейдем ко второй части — формализации задачи.
Мой индивидуальный вариант — пятый. Для программной реализации и математического анализа мне были выданы два векторных поля: поле «а» и поле «b».
Поле «а» задано функцией $y^2 i + z^2 j + x^2 k$.
Рабочая область (поверхность S) представляет собой часть единичной сферы в первом октанте.

\vspace{0.5cm}
\textbf{[Слайд 8: Анализ поля «а»]} \\
В ходе аналитического расчета поля «а» было установлено следующее:
Во-первых, дивергенция поля равна нулю, а ротор не равен нулю. Это означает, что поле является соленоидальным (или трубчатым).
Во-вторых, я рассчитал поток поля через заданную полусферу двумя способами: напрямую через поверхностный интеграл и по теореме Гаусса-Остроградского. Оба метода дали одинаковый результат: три пи делить на шестнадцать.
И в-третьих, циркуляция по контуру также была вычислена двумя независимыми методами (включая теорему Стокса) и составила ровно минус два.

\vspace{0.5cm}
\textbf{[Слайд 9: Анализ поля «b»]} \\
Для второго поля «b» ротор оказался равен нулевому вектору, что доказывает его потенциальность.
Я восстановил функцию потенциала по полному дифференциалу.
Зная потенциал, я вычислил работу сил поля по перемещению материальной точки из начала координат в точку M (один, один, один). Работа вычисляется как разность потенциалов и в моем случае составила ровно 3, что полностью совпало с контрольным ответом.

\vspace{0.5cm}
\textbf{[Слайд 10: Библиотека]} \\
Также хочу отметить, что для грамотного оформления документации я работал с первоисточниками — оригинальным четырехтомником стандартов ГОСТ ЕСПД в библиотеке 7-го корпуса МАИ. 

\vspace{0.5cm}
\textbf{[Слайд 11: Фото]} \\
\textit{(Показываем фото)} Здесь представлен фрагмент содержания четырехтомника, подтверждающий наличие стандартов.

\vspace{0.5cm}
\textbf{[Слайд 12: Финал]} \\
Проведенные расчеты станут математическим ядром для моей C++ программы FieldViz. На этом мой доклад окончен. Большое спасибо за внимание! Готов ответить на ваши вопросы.

\end{document}